\chapter{Следово-двойственный селектор межсекторного взаимодействия}
\label{ch:v8-trace-dual-cross-interaction-selector-gate}

\section{Какую свободу требуется устранить}

Микроскопический гамильтониан повторных взаимодействий предыдущей главы имеет
общий вид
\begin{equation}
 H_C=\sum_{a,b}C_{ab}D_b\otimes
 \bigl(|a\rangle\langle0|+|0\rangle\langle a|\bigr).
 \label{eq:v8-general-cross-interaction-coupling}
\end{equation}
Калибровочно-ковариантные вещественные матрицы $C$ образуют восьмимерный
коммутант. Однако редуцированный GKSL-генератор зависит не от самого $C$,
а от положительной скоростной метрики
\begin{equation}
 R=C^TC.
 \label{eq:v8-cross-coupling-rate-metric}
\end{equation}
Её калибровочный коммутантная свобода четырёхмерна. Проверим, задаёт ли уже
построенный полевой суперслед одну каноническую точку этого конуса.

\section{Точная метрика полного межсекторного модуля}

Межсекторные амплитуды образуют комплексную матрицу
\begin{equation}
 B\in M_{2\times3}(\mathbb C),
\end{equation}
где две строки соответствуют типам $QLYR$ и $XLdR$, а три столбца ---
цветовым направлениям. В полевой суперсвязности вариация конечного
оператора содержит как линейное ребро $B$, так и производную
$\Lambda^2B$. На двенадцати вещественных матричных единиц направлениях точный
SymPy-аудит без Float даёт
\begin{equation}
 K_B(\delta B_1,\delta B_2)
 =\Tr(\delta\mathcal D_B^{(1)}\delta\mathcal D_B^{(2)}),
 \qquad
 [K_B]=3I_{12}.
 \label{eq:v8-exact-cross-field-trace-metric}
\end{equation}
Следовательно, метрически двойственный модуль имеет
\begin{equation}
 K_B^{-1}=\frac13I_{12}.
 \label{eq:v8-trace-dual-cross-rate-metric}
\end{equation}

Если скачковая часть минимальной среды является именно двойственным модулем к
межсекторным амплитудам относительно этого же следа, каноническая кооценка-
свёртка задаётся
\begin{equation}
 C_{\rm tr}=K_B^{-1/2}=\frac1{\sqrt3}I_{12},
 \qquad
 C_{\rm tr}^TC_{\rm tr}=\frac13I_{12}.
 \label{eq:v8-trace-dual-canonical-coupling}
\end{equation}

\section{Ортогональная свобода является калибровкой среды}

Все вещественные матрицы связи с той же скоростной метрикой имеют вид
\begin{equation}
 C=\frac1{\sqrt3}O,
 \qquad O^TO=I_{12}.
 \label{eq:v8-cross-coupling-orthogonal-factor}
\end{equation}
Левое умножение на $O$ является заменой ортонормированного базиса
двенадцатимерной скачковой части среды, оставляющей вакуум $|0\rangle$
неподвижным. Поэтому соответствующие гамильтониан унитарно эквивалентны на
среде и дают один редуцированный канал. Восьмимерная микроскопическая
свобода не исчезает как запись оператора, но после факторизации
ненаблюдаемой калибровки репера одна скалярная скоростная метрика остаётся
единственной.

Короткий унитарный шаг имеет точную касательную
\begin{equation}
 \Psi_h(X)
 =X+\frac h3\mathcal L_{q\ell}(X)+O(h^2).
 \label{eq:v8-trace-dual-cross-gksl-tangent}
\end{equation}
Равенство генераторов проверено на всех $441$ матричных единиц полного
$M_{21}(\mathbb C)$.

\section{Совместимость с полярной геометрией}

Изотропная мобильность не уничтожает ранее полученную полярную ось. Для
точного межсекторный гессиана $H_{qX}$ выполнено
\begin{equation}
 K_B^{-1}H_{qX}=\frac13H_{qX},
\end{equation}
поэтому собственные направления и угол $55.45092\ldots^\circ$ сохраняются.
Следовая метрика определяет способ движения по уже существующему
полярному рельефу, а не заменяет сам рельеф.

\section{Граница физического статуса}

Положительный результат условен относительно одного нового, но явного
принципа: минимальная среда берётся как метрически двойственный
межсекторный модуль относительно того же конечного суперследа, который задаёт
кинетику поля. Калибровочная симметрия одна этот принцип не вынуждает.

Кроме того, замена $C_{\rm tr}\mapsto gC_{\rm tr}$ сохраняет форму
метрики и меняет только общий темп $u\mapsto g^2u$. Следовательно,
размерная единица времени, автономная поставка свежих вспомогательных систем и веса
связывающих и калибровочных семейств этим гейтом не выводятся.

\begin{theorem}[Условная единственность метрика межсекторных скоростей метрики]
На полном двенадцатимерном межсекторном модуле конечный полевой суперслед задаёт
точную метрику $K_B=3I_{12}$. При отождествлении минимальной среды с
метрически двойственным модулем каноническая свёртка взаимодействия имеет
$C^TC=I_{12}/3$. Все её квадратные корни отличаются ортогональным
переобозначением среды, поэтому редуцированный межсекторный генератор единствен с
точностью до общего масштаба времени.
\end{theorem}

\begin{nogocriterion}[Запрет переноса на полный шумовой сектор]
Нельзя объявлять этим результатом единственную метрику всех шести семейств
QMS: доказательство относится только к калибровочно-замкнутому межсекторному модулю
$QLYR\oplus XLdR$. Нельзя также считать принцип метрически двойственной
среды следствием одной калибровочной симметрии или превращать коэффициент $1/3$ в
размерную физическую скорость без независимой калибровки времени.
\end{nogocriterion}

\begin{protocol}
\begin{enumerate}
 \item межсекторный модуль: \textbf{$M_{2\times3}(\mathbb C)$, $12$ вещественных измерений};
 \item полевая следовая метрика: \textbf{$3I_{12}$, точно};
 \item двойственная скоростная метрика: \textbf{$I_{12}/3$};
 \item каноническая матрица связи: \textbf{$I_{12}/\sqrt3$};
 \item полный калибровочный коммутант $C$: \textbf{$8$ вещественных измерений};
 \item калибровочный коммутант скоростных метрик: \textbf{$4$ вещественных измерений};
 \item ортогональная свобода при фиксированном $C^TC$: \textbf{репер среды};
 \item редуцированный генератор: \textbf{$\mathcal L_{q\ell}/3$};
 \item проверенных матричных единиц: \textbf{$441$};
 \item полярная межсекторная ось: \textbf{сохранена};
 \item единственная форма метрика межсекторных скоростей метрики: \textbf{условно получена};
 \item общий масштаб времени: \textbf{не выведен};
 \item полный шестисемейный QMS: \textbf{не закрыт};
 \item следующий тест: \textbf{происхождение метрически двойственной среды
       из одного родительского действия}.
\end{enumerate}
\end{protocol}

Машинный сертификат сохранён в
\path{s2t_v8_trace_dual_cross_interaction_selector_gate_results.json}.

\begin{thebibliography}{9}
\bibitem{v8-trace-dual-carlen-maas}
E.~A.~Carlen, J.~Maas,
\emph{Gradient flow and entropy inequalities for quantum Markov semigroups
with detailed balance}, J. Funct. Anal. \textbf{273} (2017), 1810--1869;
arXiv:1609.01254.
\bibitem{v8-trace-dual-attal-pautrat}
S.~Attal, Y.~Pautrat,
\emph{From repeated to continuous quantum interactions},
Annales Henri Poincar\'e \textbf{7} (2006), 59--104;
arXiv:math-ph/0311002.
\end{thebibliography}